% 10_sop_writing.tex — LCS Workshop on writing a great SoP + Writing an Effective SoP

\section{LCS Workshop on Writing a Great SoP Statement of Purpose with Jaime Potter}
\label{sec:sop_workshop}

In organizational behavior expert Jamie Potter offers structured guidance on crafting a compelling Statement of Purpose (SOP) for PhD applications.
The session takes the form of a Q\&A, where Potter emphasizes that candidates should not view the SOP as a general application to an institution, but rather as a targeted proposal to specific faculty members.
As Potter puts it, ``you apply to a professor who is focused on a specific research area.''
A successful SOP communicates a clear understanding of a faculty member's research and proposes a direction for future collaboration.
More than half of the statement should focus on the applicant's knowledge of the professor's work and their potential to contribute.
Prior experience should be discussed only to the extent that it substantiates this contribution.

\subsection{Focusing on Faculty Research}

Candidates are advised to identify two to three faculty members whose work aligns with their own research goals.
Potter encourages applicants to study a professor's body of work and focus in particular on the ``future directions'' sections of recent papers.
The aim is to propose a credible continuation or expansion of that line of research.
As Potter puts it, ``Here is where I think we should study together over the next 5 years, and here is how I can help you.''
This proactive framing underscores the applicant's preparation and initiative.
The strongest SOPs identify synergies among multiple faculty members, since focusing on only one professor can introduce risk if that professor departs.
This triangulation also anticipates the composition of the applicant's eventual dissertation committee.

\subsection{Structure and Formality}

The SOP should be written in a formal tone, but not in a way that mimics technical reports.
Potter discourages the use of emojis, acronyms, or casual language, as the goal is to convey sincere interest in what may be the faculty member's ``life's work.''
Clear, authentic engagement with the research is prioritized over verbosity or flourish.
Potter identifies the ``gold standard'' SOP as one that demonstrates how a candidate's potential contributions lie at the intersection of two to three professors' research programs.

\subsection{Differentiating from Undergraduate Applications}

Unlike undergraduate personal statements, the SOP should not focus on academic metrics such as grades or test scores.
Once a candidate is above a baseline of competence, ``skills > grades.''
The emphasis shifts toward alignment with the faculty's research program and evidence of scholarly maturity.

\subsection{Masters vs.\ PhD Statements}

The intended degree influences SOP content.
For career-oriented master's programs, it is acceptable to emphasize coursework or broader career goals.
In contrast, research-focused master's programs and PhD tracks require SOPs that articulate research interests, demonstrate familiarity with faculty work, and indicate long-term potential for contribution.
A brief mention of one or two faculty members is sufficient in master's SOPs, but depth and specificity are expected in PhD applications.

\subsection{On Past Experience and Skill Relevance}

Job experience should only be included when it provides evidence of relevant skills or illustrates the origin of a research interest.
Potter advises against listing unrelated job responsibilities and instead encourages applicants to ``create a narrative'' that motivates their research direction.
If past academic performance was weak, candidates should not pursue a master's solely to raise grades; instead, they can demonstrate renewed competency through strong GRE scores or GRE subject tests.
To gain research experience, Potter suggests directly emailing faculty to ask, ``Is there any help I can give you?'', whether paid or unpaid, during evenings or weekends.
Importantly, demonstrating awareness of a professor's imminent research trajectory and clearly presenting oneself as the person who can carry it forward is critical---particularly because ``senior professors are very, very busy.''
Once baseline competence is met, the most salient factor is the synergy between the applicant's interests and the advisor's direction.

\subsection{Length and Supporting Materials}

SOPs should typically be limited to two single-spaced pages (approximately 1000 words).
When no page limit is specified, this length remains a reasonable target.
If applicants find themselves exceeding this, it is likely due to covering too many professors or including excessive personal history.
That said, a full page dedicated to one professor's research may be appropriate if it is not redundant; a range of one paragraph to one page per faculty member is acceptable.
Applicants may reference supporting materials such as blog posts or research documents via informal links.
As Potter suggests, a good summary statement is: ``Here is my story that has led me to these research interests with this professor, and here is how I can help you accomplish it because of my skills and experience in xyz.''

\subsection{Balancing Research Vision and Realism}

Long-term research visions are encouraged, though applicants should acknowledge that directions may evolve.
In countries such as the U.S., PhD programs span roughly five years with the first two devoted to coursework, while in countries like the UK, the PhD typically follows a two-year research master's and focuses solely on research over three years.
Consequently, SOPs in the UK may resemble formal research proposals, while U.S.\ SOPs often prioritize collaborative fit and training potential.
Potter advises applicants to signal that ``they know my work and they are skilled,'' while also showing enthusiasm for the university community more broadly.
The goal, ultimately, is to contribute to the body of academic literature---much of which is increasingly produced in industrial research labs.
When asked how to be ambitious but realistic, Potter replies that a candidate might say: ``Maybe I publish 2--3 papers, then I find my own independent direction.''
Even if a candidate's actual PhD work diverges from the SOP, the statement must still project a coherent research identity and credible trajectory.
As Potter noted, his own SOP had ``zero'' overlap with what he ended up pursuing---but it still demonstrated research potential.

\subsection{Complementing Letters of Recommendation}

Overlap between letters of recommendation and the SOP is not only acceptable but beneficial.
As Potter clarifies, redundancy in this case acts as reinforcement: ``they know what the professors are curious about and need to know about you.''
Sharing a draft SOP with recommenders can help ensure consistency in tone and message.

\subsection{Final Considerations}

Applicants should avoid unverifiable claims such as ``I am hardworking.''
Instead, desirable qualities must be demonstrated through concrete actions and examples.
The most effective SOPs blend a clear research agenda with authentic motivation and evidence.
Above all, the SOP should communicate that the applicant is prepared to contribute meaningfully to a specific scholarly community---by showing, not telling, how their interests and experience align with the research culture of the program.

\section{Writing an Effective Statement of Purpose}
\label{sec:effective_sop}
